\chapter{Introduzione alla Fisica dello Stato Solido}

\begin{nota}[Informazioni Generali]
\textbf{Data:} 29 settembre 2020\\
\textbf{Docente:} Prof. Sergio Caprara\\
\textbf{Testo di riferimento:} N.W. Ashcroft, N.D. Mermin, \textit{Solid State Physics}\\
\end{nota}

\vspace{0.5cm}

\section{Panoramica del corso}

Il corso costituisce un'introduzione alla fisica dello stato solido a livello elementare. Gli argomenti trattati sono i seguenti:

\begin{enumerate}
    \item \textbf{Struttura cristallina} - Definizione di solido cristallino e strumenti matematici per descrivere i cristalli.
    \item \textbf{Diffrazione di raggi X} - Teoria dello scattering di raggi X da parte di un cristallo, come metodo sperimentale per lo studio della struttura cristallina.
    \item \textbf{Vibrazioni reticolari (fononi)} - Descrizione delle oscillazioni degli atomi attorno alle posizioni di equilibrio e loro contributo a grandezze fisiche quali, ad esempio, il calore specifico del solido.
    \item \textbf{Proprietà elettroniche} - Trattazione della parte elettronica: come gli elettroni contribuiscono alle proprietà fisiche di un solido. In particolare, come distinguere un metallo da un isolante. Questa distinzione è una manifestazione della fisica quantistica anche a temperatura ambiente e al di sopra di essa: a causa del principio di esclusione di Pauli, gli elettroni devono disporsi nella struttura a bande del sistema. A seconda di quante bande sono occupate e del grado di riempimento di ciascuna, il sistema risulterà un metallo o un isolante.
    \item \textbf{Semiconduttori} - Studio della fisica dei semiconduttori come parte finale del corso.
\end{enumerate}

\section{Potenziale interatomico e fasi della materia}

\subsection*{Interazione tra due atomi}

A ciascun atomo possiamo attribuire un'energia cinetica $K$, che dipende dalla sua velocità, e un'energia di interazione $U$, che dipende dalla distanza tra gli atomi.

L'energia di interazione tra due atomi in funzione della loro distanza $R$ ha tipicamente il seguente andamento:

\begin{itemize}
    \item \textbf{Grandi distanze} ($R \to \infty$): $U(R) \to 0$ (per convenzione). Gli atomi non esercitano forze l'uno sull'altro.
    \item \textbf{Distanza intermedia}: esiste una \textbf{distanza di equilibrio} $R_0$ corrispondente al minimo del potenziale, di profondità $U_0 < 0$.
    \item \textbf{Piccole distanze} ($R \to 0$): forte \textbf{repulsione} dovuta alla sovrapposizione delle shell elettroniche esterne. Non è possibile avvicinare eccessivamente due atomi.
\end{itemize}



Il valore $|U_0|$ (profondità della buca di potenziale) fornisce una stima dell'energia di interazione tipica tra due atomi.

\subsection*{Parametri termodinamici e fasi}

Le variabili che controllano il comportamento macroscopico di un insieme di atomi sono la \textbf{temperatura} $T$ e la \textbf{pressione} $P$ (o equivalentemente la \textbf{densità} $\rho$).

\subsubsection{Fase gassosa: $K \gg |U_0|$}

A temperature molto elevate o densità molto basse, l'energia cinetica degli atomi è molto maggiore dell'energia di interazione tipica. In queste condizioni, nessun atomo può essere intrappolato nella buca di potenziale, poiché l'energia cinetica è sufficiente a farlo sfuggire. Il sistema si trova in \textbf{fase gassosa}.

La fase gassosa è la fase più semplice della materia. Nel limite del \textbf{gas perfetto} (temperatura molto alta o densità molto bassa), è possibile calcolare esattamente l'equazione di stato. Questo è l'unico sistema per il quale si conosce la soluzione esatta in meccanica statistica.

\subsubsection{Condensazione: $K \sim |U_0|$}

Riducendo la temperatura o aumentando la densità (ad esempio riducendo il volume o aumentando la pressione), si raggiunge un punto in cui le due grandezze sono comparabili. Si verifica allora il fenomeno della \textbf{condensazione}, che dà luogo a una \textbf{fase condensata}. La fisica della materia condensata si occupa appunto delle fasi condensate (non gassose).

\subsubsection{Fase liquida: $K \sim |U_0|$}

Un \textbf{liquido} è un sistema con le seguenti proprietà:
\begin{itemize}
    \item \textbf{Comprimibilità molto bassa}: è difficile comprimere un liquido (a differenza di un gas).
    \item \textbf{Volume ben definito}.
    \item \textbf{Gli atomi possono muoversi} gli uni rispetto agli altri, poiché l'energia cinetica è comparabile con l'energia potenziale.
    \item Esiste una \textbf{distanza media caratteristica} tra gli atomi (dell'ordine di $R_0$), ma su grande scala gli atomi occupano posizioni casuali.
    \item Il liquido può \textbf{fluire} perché gli atomi sono relativamente liberi di muoversi.
\end{itemize}

\subsubsection{Fase solida: $K \ll |U_0|$}

A temperature molto basse o pressioni molto elevate, l'energia cinetica è molto minore dell'energia potenziale. In un \textbf{solido}, ogni atomo è intrappolato nella buca di potenziale e la sua energia cinetica consente soltanto \textbf{piccole oscillazioni} attorno alla posizione di equilibrio, senza possibilità di sfuggire.

Esistono essenzialmente due tipi di solido:
\begin{itemize}
    \item \textbf{Solido amorfo} (es. vetro, zolfo amorfo): gli atomi sono congelati in posizioni fisse, il sistema non può fluire, ma le posizioni degli atomi sono \textbf{casuali} su larga scala. Non c'è ordine a lungo raggio.
    \item \textbf{Solido cristallino}: gli atomi sono non solo congelati in posizioni fisse, ma queste posizioni formano una \textbf{distribuzione ordinata} nello spazio, un pattern ben definito che costituisce un \textbf{cristallo}.
\end{itemize}

La fase cristallina è una fase a \textbf{entropia molto bassa}, poiché gli atomi sono ordinati nello spazio in una struttura ben definita.

\subsection*{Caso speciale: l'elio come liquido quantistico}

Quasi tutti gli elementi sono solidi a temperatura sufficientemente bassa, \textbf{tranne l'elio}. L'elio resta liquido fino allo zero assoluto per due ragioni:
\begin{enumerate}
    \item È un elemento molto \textbf{leggero}, quindi la sua energia cinetica può essere significativa anche a basse temperature (effetti quantistici).
    \item È un \textbf{gas nobile}, quindi l'attrazione tra atomi di elio è \textbf{estremamente debole} (il potenziale $|U_0|$ è molto piccolo).
\end{enumerate}

L'elio è pertanto un \textbf{liquido quantistico}: gli effetti quantistici (in particolare, l'incertezza nella posizione dei nuclei, che è molto maggiore della distanza media tra essi) diventano importanti prima che il sistema possa cristallizzare.

\begin{nota}[Approssimazione del corso]
Nel corso tratteremo tutti i sistemi come solidi cristallini a temperatura zero, ignorando questi effetti quantistici peculiari dell'elio.
\end{nota}

\section{Reticoli di Bravais}

\subsection*{Definizione formale}

Si consideri uno spazio tridimensionale. Si assegnino tre vettori $\mathbf{a}_1$, $\mathbf{a}_2$, $\mathbf{a}_3$ \textbf{linearmente indipendenti} (non coplanari).

\begin{definizione}{Reticolo di Bravais}
Un reticolo di Bravais è l'insieme di tutti i vettori $\mathbf{R}$ della forma:
\begin{equation}
\mathbf{R} = n_1 \mathbf{a}_1 + n_2 \mathbf{a}_2 + n_3 \mathbf{a}_3, \qquad n_1, n_2, n_3 \in \mathbb{Z}
\end{equation}
\end{definizione}

I vettori $\mathbf{a}_1$, $\mathbf{a}_2$, $\mathbf{a}_3$ sono detti \textbf{vettori primitivi}. Un reticolo di Bravais è una struttura \textbf{infinita} in tutte le direzioni e \textbf{periodica}.


Un cristallo reale è finito e possiede bordi. Nessun cristallo reale è quindi, rigorosamente, un reticolo di Bravais. Tuttavia, lontano dai bordi, il modello del reticolo di Bravais costituisce un'ottima descrizione della struttura interna di un cristallo reale. Si tratta di un \textbf{modello matematico} che approssima il mondo fisico; occorre sempre tenere presente la differenza tra modello e realtà, e adattare la descrizione quando le assunzioni del modello vengono meno.


\subsection*{Proprietà fondamentali}

\subsubsection{Periodicità}
Il reticolo appare identico da qualunque punto reticolare si osservi. Spostandosi da un punto del reticolo a un altro, la vista dell'intorno è esattamente la stessa, poiché il reticolo è infinito in tutte le direzioni.

\begin{nota}[Simmetria Traslazionale]
\textbf{Tutti i punti di un reticolo di Bravais sono equivalenti.}
\end{nota}

\subsubsection{Simmetria di inversione}
Se $\mathbf{R}$ appartiene al reticolo, allora anche $-\mathbf{R}$ appartiene allo stesso reticolo.

\textbf{Dimostrazione:} se $n_1, n_2, n_3$ sono interi che definiscono il punto $\mathbf{R}$, allora $-n_1, -n_2, -n_3$ sono anch'essi interi e definiscono il punto $-\mathbf{R}$.

Quindi, \textbf{per definizione}, tutti i reticoli di Bravais godono di:
\begin{itemize}
    \item \textbf{Simmetria traslazionale} (periodicità)
    \item \textbf{Simmetria di inversione}
\end{itemize}

\subsection*{Reticolo di Bravais in una dimensione}

In una dimensione, esiste \textbf{un solo} reticolo di Bravais possibile: tutti gli atomi disposti sulla retta a distanza costante $a$ l'uno dall'altro.

\vspace{0.5cm}
\begin{center}
\begin{tikzpicture}
\foreach \x in {0, 2, 4, 6, 8} {
    \filldraw[colorA] (\x,0) circle (4pt);
}
\draw[<->, thick, color=colorD] (0.2,0.3) -- (1.8,0.3) node[midway, above] {$a$};
\draw[<->, thick, color=colorD] (2.2,0.3) -- (3.8,0.3) node[midway, above] {$a$};
\draw[<->, thick, color=colorD] (4.2,0.3) -- (5.8,0.3) node[midway, above] {$a$};
\draw[<->, thick, color=colorD] (6.2,0.3) -- (7.8,0.3) node[midway, above] {$a$};
\node[scale=1.5] at (9,0) {$\cdots$};
\node[scale=1.5] at (-1,0) {$\cdots$};
\end{tikzpicture}
\end{center}
\vspace{0.5cm}

Questo è un reticolo monodimensionale con passo $a$, ed è l'unica disposizione ordinata possibile (a meno di permutazioni di atomi indistinguibili).

I reticoli cristallini furono classificati dal fisico \textbf{Auguste Bravais}, da cui il nome "reticolo di Bravais".

\section{Strutture non-Bravais}

\subsection*{Il reticolo a nido d'ape (honeycomb)}
Non tutte le strutture periodiche sono reticoli di Bravais. Consideriamo la struttura a \textbf{nido d'ape} (honeycomb): una rete di esagoni regolari senza punti centrali.

Questa struttura \textbf{non è} un reticolo di Bravais, perché non tutti i punti sono equivalenti. Se ci si pone su un punto del reticolo e si osserva l'intorno, la vista dipende dalla posizione scelta:
\begin{itemize}
    \item Da un punto (es. di colore "rosa"), si vede un vicino a sinistra a $0^\circ$, uno a $+120^\circ$ e uno a $-120^\circ$.
    \item Da un punto adiacente (es. di colore "viola"), si vede un vicino a destra a $0^\circ$, uno a $+120^\circ$ a sinistra e uno a $-120^\circ$ a sinistra.
\end{itemize}
I due punti non sono equivalenti, violando la condizione fondamentale dei reticoli di Bravais.

\subsection*{Descrizione tramite reticolo con base}
Osservando attentamente, i punti dell'honeycomb possono essere suddivisi in \textbf{due classi} (es. "rosa" e "viola"):
\begin{itemize}
    \item Tutti i punti \textbf{rosa} sono equivalenti tra loro e formano un reticolo di Bravais.
    \item Tutti i punti \textbf{viola} sono equivalenti tra loro e formano un altro reticolo di Bravais (dello stesso tipo).
\end{itemize}
La struttura complessiva è quindi la \textbf{compenetrazione di due reticoli di Bravais}, traslati l'uno rispetto all'altro di un \textbf{vettore di base}. Questa compenetrazione è chiamata \textbf{reticolo con base} (o reticolo di Bravais con base).

\subsection*{Honeycomb con punti centrali}
Se si aggiunge il \textbf{punto centrale} a ciascun esagono dell'honeycomb, la struttura risultante \textbf{è} un reticolo di Bravais (reticolo esagonale/triangolare). In questo caso, tutti i punti sono davvero equivalenti.

\section{Reticoli di Bravais in due dimensioni (5 reticoli)}

In due dimensioni esistono \textbf{5 reticoli di Bravais}.

\subsubsection{Reticolo monoclino (obliquo)}
\begin{itemize}
    \item $|\mathbf{a}_1| \neq |\mathbf{a}_2|$
    \item Angolo $\varphi \neq 90^\circ$
    \item \textbf{Simmetrie}: solo traslazione e inversione (nessuna simmetria aggiuntiva).
    \item La cella è un parallelogramma generico.
\end{itemize}
Questo è il reticolo meno simmetrico in 2D.

\subsubsection{Reticolo ortorombico (rettangolare)}
\begin{itemize}
    \item $|\mathbf{a}_1| \neq |\mathbf{a}_2|$
    \item Angolo $\varphi = 90^\circ$
    \item \textbf{Simmetrie}: traslazione, inversione, \textbf{rotazione di $180^\circ$} attorno a un asse perpendicolare al piano, \textbf{simmetria speculare} (mirror) rispetto a due piani.
    \item La cella è un rettangolo.
\end{itemize}

\begin{definizione}[Distinzione tra monoclino e ortorombico]
\textbf{Distinzione tra monoclino e ortorombico:} la rotazione di $180^\circ$ è in realtà presente anche nel monoclino in 2D (equivale all'inversione in 2D). Ciò che distingue i due reticoli è la \textbf{simmetria speculare} (mirror symmetry): il rettangolare possiede piani speculari lungo le direzioni principali, mentre il monoclino (obliquo) no. In un reticolo obliquo, la riflessione rispetto alle direzioni principali non mappa il reticolo in sé stesso.
\end{definizione}

\subsubsection{Reticolo esagonale}
\begin{itemize}
    \item $|\mathbf{a}_1| = |\mathbf{a}_2|$
    \item Angolo $\varphi = 120^\circ$
    \item Il punto centrale dell'esagono si ottiene come $\mathbf{a}_1 + \mathbf{a}_2$ (partendo dall'origine).
\end{itemize}

\subsubsection{Reticolo tetragonale (quadrato)}
\begin{itemize}
    \item $|\mathbf{a}_1| = |\mathbf{a}_2|$
    \item Angolo $\varphi = 90^\circ$
    \item \textbf{Simmetrie}: gode di \textbf{rotazione di $90^\circ$} (che il rettangolare non possiede). Ruotando il reticolo di $90^\circ$, esso si sovrappone a sé stesso.
    \item Il reticolo più simmetrico in 2D.
\end{itemize}

\subsubsection{Reticolo ortorombico centrato}
\begin{itemize}
    \item $|\mathbf{a}_1| = |\mathbf{a}_2|$
    \item $\varphi \neq 90^\circ$ e $\varphi \neq 120^\circ$ (per escludere tetragonale ed esagonale)
    \item Si ottiene partendo dal reticolo ortorombico e aggiungendo un punto al centro di ciascun rettangolo.
    \item Possiede le \textbf{stesse simmetrie} dell'ortorombico, ma è un reticolo diverso.
\end{itemize}

\subsubsection{Riepilogo}

\begin{table}[h!]
\centering
\renewcommand{\arraystretch}{1.3}
\begin{tabular}{|l|c|l|l|}
\hline
\textbf{Reticolo} & \textbf{$\vert\mathbf{a}_1\vert$ vs $\vert\mathbf{a}_2\vert$} & \textbf{Angolo} & \textbf{Simmetria caratteristica} \\ \hline
Monoclino (obliquo) & $\neq$ & $\varphi$ generico & Solo inversione \\ \hline
Ortorombico (rettangolare) & $\neq$ & $90^\circ$ & Mirror symmetry \\ \hline
Esagonale & $=$ & $120^\circ$ & Rotazione $60^\circ$ \\ \hline
Tetragonale (quadrato) & $=$ & $90^\circ$ & Rotazione $90^\circ$ \\ \hline
Ortorombico centrato & $=$ & $\varphi \neq 90^\circ, 120^\circ$ & Mirror symmetry \\ \hline
\end{tabular}
\end{table}

In sintesi: \textbf{4 reticoli fondamentali + 1 decorazione} (l'ortorombico centrato) = \textbf{5 reticoli di Bravais in 2D}.

\section{Reticoli di Bravais in tre dimensioni (14 reticoli)}

In tre dimensioni esistono in totale \textbf{14 reticoli di Bravais}: \textbf{7 fondamentali} (con le simmetrie principali) e \textbf{7 ottenuti per decorazione} (aggiunta di punti supplementari).

\subsection*{I 7 sistemi cristallini fondamentali}

\subsubsection{1. Triclino}
\begin{itemize}
    \item $a_1 \neq a_2 \neq a_3$ (tutti i lati diversi)
    \item $\alpha \neq \beta \neq \gamma$ (tutti gli angoli diversi)
    \item Il parallelepipedo ha tutte le facce diverse (parallelogrammi diversi).
    \item \textbf{Simmetrie}: solo inversione. Nessuna simmetria speculare.
    \item È il sistema \textbf{meno simmetrico} di tutti.
    \item \textbf{Non ammette decorazioni}: 1 solo reticolo triclino.
\end{itemize}

\subsubsection{2. Monoclino}
\begin{itemize}
    \item $a_1 \neq a_2 \neq a_3$ (tutti i lati diversi)
    \item \textbf{Due angoli sono $90^\circ$}, un solo angolo è diverso da $90^\circ$.
    \item Conseguenza: due facce sono \textbf{rettangoli} e una faccia è un \textbf{parallelogramma}.
    \item \textbf{Simmetrie}: inversione + \textbf{simmetria speculare} rispetto al piano rettangolare. Più simmetrico del triclino.
    \item \textbf{Decorazione possibile}: base centrata (un punto aggiuntivo al centro della base e del top).
    \item \textbf{Totale: 2 reticoli monoclini} (semplice + base centrata).
\end{itemize}

\subsubsection{3. Ortorombico}
\begin{itemize}
    \item $a_1 \neq a_2 \neq a_3$ (tutti i lati diversi)
    \item \textbf{Tutti gli angoli sono $90^\circ$}: tutte le facce sono rettangoli (ma tutti diversi tra loro).
    \item \textbf{Simmetrie}: molte rotazioni e simmetrie speculari.
    \item \textbf{Decorazioni possibili} (3):
    \begin{itemize}
        \item \textbf{Base centrata}: un punto al centro della base (e del top).
        \item \textbf{Facce centrate} (\textit{face centered}): un punto al centro di ciascuna delle 6 facce.
        \item \textbf{Corpo centrato} (\textit{body centered}): un punto all'incrocio delle diagonali del parallelepipedo.
    \end{itemize}
    \item \textbf{Totale: 4 reticoli ortorombici} (semplice + base centrata + facce centrate + corpo centrato).
\end{itemize}

\subsubsection{4. Tetragonale}
\begin{itemize}
    \item Base quadrata: $a_1 = a_2 \neq a_3$.
    \item \textbf{Tutti gli angoli sono $90^\circ$}.
    \item È un parallelepipedo con una base quadrata e altezza diversa dal lato della base.
    \item \textbf{Decorazione possibile}: solo \textbf{corpo centrato} (\textit{body centered}).
    \item Il \textbf{tetragonale base centrata} non è un nuovo reticolo: aggiungendo un punto al centro della base quadrata, si ottiene un nuovo reticolo quadrato (ruotato di $45^\circ$ e più piccolo), che è semplicemente un altro reticolo tetragonale semplice.
    \item Analogamente, il \textbf{tetragonale facce centrate} si riduce a un reticolo già presente nella lista.
    \item \textbf{Totale: 2 reticoli tetragonali} (semplice + corpo centrato).
\end{itemize}

\subsubsection{5. Romboedrico (trigonale)}
\begin{itemize}
    \item \textbf{Tutti i lati uguali}: $a_1 = a_2 = a_3$.
    \item \textbf{Tutti gli angoli uguali} in ciascun vertice, ma diversi da $90^\circ$.
    \item Si ottiene dal cubo "schiacciandolo" lungo una diagonale. Le facce sono \textbf{rombi}.
    \item \textbf{Non ammette decorazioni}: 1 solo reticolo romboedrico.
\end{itemize}

\subsubsection{6. Esagonale}
\begin{itemize}
    \item La cella di base è un \textbf{rombo con angolo di $120^\circ$}, formato da due triangoli equilateri.
    \item Si chiama "esagonale" perché \textbf{tre} di queste celle affiancate formano un \textbf{esagono regolare}.
    \item $a_1 = a_2 \neq a_3$, angolo alla base = $120^\circ$.
    \item \textbf{Non ammette decorazioni}: 1 solo reticolo esagonale.
\end{itemize}

\subsubsection{7. Cubico}
\begin{itemize}
    \item $a_1 = a_2 = a_3$, tutti gli angoli = $90^\circ$.
    \item Il \textbf{cubo} è l'oggetto \textbf{più simmetrico} tra tutti i sistemi cristallini in 3D.
    \item Possiede moltissime simmetrie. Ad esempio, una \textbf{rotazione di $120^\circ$} attorno alla diagonale del corpo mappa il cubo in sé stesso. Esiste un \textbf{esagono nascosto} nel cubo: la sezione del cubo lungo il piano perpendicolare alla diagonale del corpo è un esagono regolare.
    \item \textbf{Decorazioni possibili} (2):
    \begin{itemize}
        \item \textbf{Facce centrate} (FCC, \textit{face-centered cubic}): un punto al centro di ciascuna delle 6 facce.
        \item \textbf{Corpo centrato} (BCC, \textit{body-centered cubic}): un punto al centro del cubo.
    \end{itemize}
    \item Il \textbf{cubico base centrata} non è un nuovo reticolo: aggiungendo un punto al centro delle due basi di un cubo, i lati della nuova base quadrata risultano uguali, ma l'altezza diventa maggiore del lato. Si ottiene quindi un \textbf{reticolo tetragonale}, già presente nella lista.
    \item \textbf{Totale: 3 reticoli cubici} (semplice + FCC + BCC).
\end{itemize}

\subsection*{Riepilogo dei 14 reticoli di Bravais}

\begin{table}[h!]
\centering
\renewcommand{\arraystretch}{1.3}
\begin{tabular}{|l|c|c|c|}
\hline
\textbf{Sistema cristallino} & \textbf{Fondamentale} & \textbf{Decorazioni} & \textbf{Totale} \\ \hline
Triclino & 1 & 0 & \textbf{1} \\ \hline
Monoclino & 1 & 1 (base centrata) & \textbf{2} \\ \hline
Ortorombico & 1 & 3 (base c., facce c., corpo c.) & \textbf{4} \\ \hline
Tetragonale & 1 & 1 (corpo centrato) & \textbf{2} \\ \hline
Romboedrico & 1 & 0 & \textbf{1} \\ \hline
Esagonale & 1 & 0 & \textbf{1} \\ \hline
Cubico & 1 & 2 (FCC, BCC) & \textbf{3} \\ \hline
\textbf{Totale} & \textbf{7} & \textbf{7} & \textbf{14} \\ \hline
\end{tabular}
\end{table}

\section{Cella primitiva}

I vettori $\mathbf{a}_1$, $\mathbf{a}_2$, $\mathbf{a}_3$ sono detti \textbf{vettori primitivi}. Il parallelepipedo costruito su questi tre vettori come spigoli è la \textbf{cella primitiva}.

\begin{equazione}[Proprietà Fondamentale]
Ogni cella primitiva contiene esattamente \textbf{un punto reticolare}.
\end{equazione}

\textbf{Dimostrazione (esempio del cubico semplice):} la cella cubica ha 8 vertici, ciascuno dei quali è un punto del reticolo. Ogni vertice è condiviso da 8 cubi adiacenti, quindi la frazione di ciascun punto appartenente a una singola cella è $\frac{1}{8}$. Il numero totale di punti per cella è:
$$8 \times \frac{1}{8} = 1$$

\subsection*{Volume della cella primitiva}
Il volume della cella primitiva è dato da:
\begin{equazione}{Volume della cella primitiva}
\begin{equation}
V = \left| \mathbf{a}_1 \cdot (\mathbf{a}_2 \times \mathbf{a}_3) \right|
\end{equation}
\end{equazione}
Questa formula vale per tutte le celle primitive, indipendentemente dal sistema cristallino.

\textbf{Esempio:} per il cubico semplice con lato $a$, si ha $\mathbf{a}_2 \times \mathbf{a}_3$ parallelo ad $\mathbf{a}_1$ e di modulo $a^2$, quindi $V = a^3$.

\subsection*{Non unicità della scelta dei vettori primitivi}
La scelta dei vettori primitivi \textbf{non è unica}. Ad esempio, per un reticolo quadrato 2D con vettori $\mathbf{a}_1$ e $\mathbf{a}_2$, una scelta alternativa valida è $\mathbf{a}_1$ e $\mathbf{a}_2' = \mathbf{a}_1 + \mathbf{a}_2$: i vettori $\mathbf{a}_1$ e $\mathbf{a}_2'$ con coefficienti interi generano esattamente lo stesso reticolo.

Esistono \textbf{infinite} scelte possibili per i vettori primitivi. Tuttavia, \textbf{tutte le celle primitive hanno lo stesso volume}, poiché tutte devono contenere esattamente un punto reticolare (altrimenti la densità del cristallo non sarebbe univocamente definita).

La scelta più ragionevole è quella che rende evidenti le \textbf{simmetrie} del reticolo.

\section{Cella convenzionale}

La cella primitiva non sempre evidenzia le simmetrie del reticolo. Ad esempio:
\begin{itemize}
    \item La cella primitiva del reticolo \textbf{esagonale} è un rombo con angolo di $120^\circ$, che non ricorda la simmetria esagonale.
    \item La cella primitiva del \textbf{cubico a facce centrate (FCC)} è un parallelepipedo molto obliquo, che non somiglia affatto a un cubo.
\end{itemize}
Per questo motivo, si introduce la \textbf{cella convenzionale}: una cella (non necessariamente primitiva) che mette in evidenza le simmetrie del sistema.

\subsection*{Esempio: cubico a facce centrate (FCC)}
La cella convenzionale dell'FCC è il \textbf{cubo} con i punti aggiuntivi al centro di ciascuna faccia. Questa cella contiene \textbf{4 punti reticolari}:
\begin{itemize}
    \item 8 vertici $\times \frac{1}{8} = 1$
    \item 6 facce $\times \frac{1}{2} = 3$
    \item Totale: $1 + 3 = 4$
\end{itemize}

I \textbf{vettori primitivi} dell'FCC sono:
\begin{equazione}{Vettori primitivi FCC}
\begin{equation}
\mathbf{a}_1 = \frac{a}{2}(1, 1, 0), \qquad \mathbf{a}_2 = \frac{a}{2}(1, 0, 1), \qquad \mathbf{a}_3 = \frac{a}{2}(0, 1, 1)
\end{equation}
\end{equazione}
dove $a$ è il lato del cubo.

Il volume della cella primitiva è:
\begin{equazione}{Volume cella primitiva FCC}
\begin{equation}
V_{\text{primitiva}} = \frac{1}{4} \, V_{\text{cubo}} = \frac{a^3}{4}
\end{equation}
\end{equazione}
poiché la cella convenzionale contiene 4 punti e la primitiva ne contiene 1.

\section{Cella di Wigner-Seitz}

Esiste una scelta della cella che è contemporaneamente:
\begin{itemize}
    \item \textbf{Primitiva}: contiene esattamente 1 punto reticolare.
    \item \textbf{Convenzionale}: mostra tutte le simmetrie del reticolo.
\end{itemize}

\begin{definizione}{Cella di Wigner-Seitz}
Questa cella è la \textbf{cella di Wigner-Seitz}. Il punto reticolare si trova esattamente al centro della cella.
\end{definizione}


\subsection*{Esempi}

\subsubsection{FCC (cubico a facce centrate)}
La cella di Wigner-Seitz dell'FCC è un poliedro (non un cubo) che possiede tuttavia \textbf{tutte le simmetrie del cubo}:
\begin{itemize}
    \item \textbf{Rotazione di $90^\circ$} (simmetria \textit{fourfold}) attorno agli assi passanti per i vertici del poliedro corrispondenti alle facce del cubo.
    \item \textbf{Rotazione di $120^\circ$} (simmetria \textit{threefold}) attorno agli assi passanti per i vertici del poliedro corrispondenti alle diagonali del cubo. L'esagono nascosto nel cubo è visibile nella cella di Wigner-Seitz dell'FCC.
\end{itemize}

\subsubsection{BCC (cubico a corpo centrato)}
La cella di Wigner-Seitz del BCC è un poliedro formato da \textbf{esagoni e quadrati}, anch'esso con tutte le simmetrie del cubo.

\subsection*{Riempimento dello spazio}
Le celle di Wigner-Seitz, sebbene di forma insolita, possono \textbf{riempire completamente lo spazio} senza lasciare vuoti né sovrapposizioni, esattamente come dei mattoncini (blocchi Lego). Sono quindi celle primitive a tutti gli effetti, con il vantaggio aggiuntivo di preservare tutte le simmetrie del sistema cristallino di appartenenza.
